\section{Dimension of the exterior-blind store}
\label{sec:h-bhq}

Lemma~\ref{lem:fiber} is the local count: one generic merge hides a complex
line, real dimension \(2\), dropping to \(1\) only in the all-zero merged state.
Proposition~\ref{prop:vol} is the global count under a boundary-sized imprint:
\[
D_Q\ge 2(N_Q-b),
\]
volume-leading when \(N_Q=n_R-1\).
If \(Q\) does not enter \(H\), the Jacobian vanishes and the kernel is the
full \(2(n_R-1)\), up to all-zero exceptions.

You could compare this with \(\sum h_2(P_-)\) from Section~\ref{sec:bond}.
The comparison is an inequality of different quantities.
An area law for one does not force an area law for the other.
The factor \(2\) is \(\dim_\R\C\).
I1a's one bit per saturated leg is a third normalization, available only
after an edge set of cardinality \(k\) and a mean \(\bar h\) have been chosen.
Theorem~\ref{thm:disjoint} fixes both the mean and the variance of a Haar
matching, and neither equals \(D_Q\).
Proposition~\ref{prop:circle} is the geometric reason:
\(h_2\) uses the radius of the fibre circle, while \(D_Q\) counts the plane.
